% sec-10-references-backmatter.tex - Back matter and references.  FULL stage.
% Follows the pattern of
% funding/capitalization-plan/final-capital/publication/sections/
% sec-11-references-backmatter.tex: abbreviations, data and code availability,
% author contributions and conflicts, citation, then the two-column reference
% list.
%
% ONE DELIBERATE DEPARTURE FROM THE PARENT.  The parent numbers its
% abbreviations block as a table.  This paper's table budget is exactly
% twenty-five, all twenty-five carry an argument, and all twenty-five are
% indexed in Table 3.  The abbreviations block is a reference aid rather than an
% argument, so it is set at the body measure in the same idiom as every other
% table but is left unnumbered, and the count of numbered tables stays at
% twenty-five.
%
% BIBLIOGRAPHY STYLE.  unsrturl, not unsrt.  unsrt drops the doi and url fields
% that references.bib carries on every entry, which is why a reference list
% built with it prints neither a DOI string nor a clickable target.  unsrturl is
% unsrt entry for entry with those two fields restored, and trialstyle.sty marks
% the printed DOIs and URLs in the accent color.  \UrlBreaks is re-asserted
% after url and hyperref so no reference line runs past the right margin.
\section*{Back Matter}
\phantomsection\addcontentsline{toc}{section}{Back Matter}

\bmhead{Abbreviations}

Every abbreviation this paper uses is expanded below, in two column pairs at the
body measure. Trial-specific terms follow the definitions in the published Phase
1 protocol rather than being redefined here.

\begin{apptable}
\begin{tabularx}{\textwidth}{>{\raggedright\arraybackslash}p{1.5cm} >{\raggedright\arraybackslash}X >{\raggedright\arraybackslash}p{1.5cm} >{\raggedright\arraybackslash}X}
\headrow \thead{Short} & \thead{Expansion} & \thead{Short} & \thead{Expansion}\\
\midrule
AE & Adverse event & IRB & Institutional review board \\
BICR & Blinded independent central review & ISGPS & \scriptsize{International Study Group of Pancreatic Surgery} \\
CCDI & Childhood Cancer Data Initiative & LLM & Large language model \\
CFR & Code of Federal Regulations & MPR & Major pathologic response \\
CMC & Chemistry, manufacturing and control & OS & Overall survival \\
CRO & Contract research organization & PDAC & Pancreatic ductal adenocarcinoma \\
ctDNA & Circulating tumor DNA & PFS & Progression-free survival \\
DLT & Dose-limiting toxicity & QSP & Quantitative systems pharmacology \\
DOI & Digital object identifier & RP2D & Recommended Phase 2 dose \\
EO & Executive Order & SAE & Serious adverse event \\
F and A & Facilities and administrative cost & SBIR & Small Business Innovation Research \\
FDA & Food and Drug Administration & STTR & Small Business Technology Transfer \\
FTE & Full-time equivalent & USL & Unified Safety Level \\
HHS & Department of Health and Human Services & VVUQ & \scriptsize{Verification, validation and uncertainty quantification} \\
ICH & \scriptsize{International Council for Harmonization} & 3+3 & Standard 3 plus 3 dose escalation \\
CMS & Centers for Medicare and Medicaid Services & ASME & \scriptsize{American Society of Mechanical Engineers} \\
HPB & Hepato-pancreato-biliary & RECIST & \scriptsize{Response Evaluation Criteria in Solid Tumors} \\
MAHA & Make America Healthy Again & NCI & National Cancer Institute \\
IDE & Investigational device exemption & & \\
IND & Investigational new drug & & \\
\end{tabularx}
\end{apptable}

\bmhead{Data and Code Availability}

{\sloppy Every work cited to a digital object identifier is openly deposited.
The repository is at
\url{https://github.com/kevinkawchak/physical-ai-oncology-trials},
version v4.6.0 \cite{repo2026}.\par} The directory
\appfile{new-trial-system} contains the eight-stage sub-prompt schedule, the
twenty-five figure specifications across five machine-readable diagram
platforms, the three build stages with an Overleaf bundle for each, the master
prompt filed verbatim together with the full build output, and the TikZ source
of every figure printed here. No raster image is used anywhere in this paper.

Four author-final directories were read directly and are the source of every
quotation in \S3 through \S6:
\appfile{trial-ind/final-ind/publication},
\appfile{trial-protocol/final-protocol/publication},
\appfile{trial-phase-2/final-protocol/publication/author}, and
\appfile{funding/capitalization-plan/final-capital/publication}. Four input
archives were read directly and are the source of every quotation in \S5 and
\S7: the three legislation archives and the AI peer review archive under
\appfile{new-trial-system/inputs}. Two further funding directories were read for
\S6: \appfile{funding/pdac-funding-applications/final-apply/publication} and
\appfile{funding/RFA-RM-27-001-v2}.

\bmhead{Author Contributions and Conflicts}

Kevin Kawchak is the sole author. He is the sole officer, sole director and sole
owner of ChemicalQDevice, which is the sponsor described throughout the
documents this paper reports on, and he holds no equity in Revolution Medicines
or in any robotic surgery vendor. He is not, and under the plan reported in
\S6 will not become, a clinical investigator on the trial. No sponsor, contract
research organization, institutional review board, regulator, or medical society
reviewed this paper or authorized any statement in it. Nothing reported here is
an approved protocol, an active IND, or an enacted statute. Claude Code Opus 5
adapted this work; ChatGPT Codex and Google Gemini served as independent
reviewers of the work reported in \S7, and none of the three is an applicant,
investigator, sponsor, regulator, clinician, or decision-maker.

\bmhead{Citation}

Kawchak K. Pancreatic Cancer LLM Clinical Trial System: From IND to Protocol to
Legislation, Funding, and AI Peer Review. Zenodo. 2026;
\dlink{10.5281/zenodo.xxxxxxxx}.

\vspace{0.2em}
\apprule{0.5pt}
\vspace{-0.1em}

\bmhead{References}
\vspace{-0.35em}
\begin{apprefs}
\bibliographystyle{unsrturl}
\bibliography{references}
\end{apprefs}
